\documentclass{situationist-preprint}

\title{Definite Descriptions and Context}
\msid{LNG-19-b}
\msfamily{LNG-19}
\msdiscipline{Philosophy of Language}
\mstype{Companion manuscript}
\mskeywords{definite descriptions; incomplete descriptions; domain restriction; pragmatic enrichment; Russell; context-sensitivity}

\begin{document}
\maketitle

\begin{abstract}
LNG-19-a argues that definite descriptions are not context-sensitive: an incomplete description such as \enquote{the table} has its complete Russellian content in every context, and the restriction that makes it apt belongs to what the speaker communicates. We develop the case for incomplete descriptions. Kripke's test of a language stipulated to be Russellian shows that the use of incomplete descriptions is no evidence for semantic restriction; the indeterminacy of completions, often taken to tell against the Russellian, tells against the semantic restriction theorist instead; and the strongest arguments for semantic domain restriction, from restrictions that vary within a sentence and from restrictions inside attitude reports, are answered by the pragmatic account.
\end{abstract}

\section{The problem}

\citet{Strawson1950} observed that \enquote{the table is covered with books} is used truly, although there is more than one table in the world. On Russell's theory the sentence says that there is exactly one table and that it is covered with books, which is false. LNG-19-a holds that this is what the sentence says (Lemma 3.2 of LNG-19-a) and that the speaker communicates a restricted proposition, about the table in the room, by pragmatic enrichment (Principle P2 of LNG-19-a). Its rival holds that the domain of the description, like that of a quantifier, is restricted in the semantic content, by a contextually valued variable or by an elliptical completion \citep{StanleySzabo2000}. By Definition 2.2 of LNG-19-a, the rival makes \enquote{the} context-sensitive. We give three further considerations against it.

\section{Three considerations}

\paragraph{Kripke's test.} \citet{Kripke1977} proposed a test for whether a phenomenon requires a semantic ambiguity: imagine a language in which the unambiguous semantics is stipulated, and ask whether speakers of it would exhibit the phenomenon anyway. Imagine, then, a language stipulated to be Russellian, whose \enquote{the $F$} says that exactly one thing in the universe is $F$. Its speakers would still say \enquote{the table is covered with books} in a room with one table. They would do so because saying \enquote{the table in this room} is more cumbersome, and because the hearer can recover the intended restriction with no more effort than in the case of quantifiers, by the reasoning that Principle P2 of LNG-19-a describes. Since the phenomenon would arise in the stipulated language, it is no evidence that our language differs from it.

\paragraph{Indeterminacy.} \citet{Wettstein1981} argued that no completion is privileged: \enquote{the table in this room}, \enquote{the table at which we are sitting}, \enquote{the table I bought yesterday} all fit, and nothing in the context selects one. He took this as an objection to the Russellian. It is rather an objection to semantic restriction. If the restriction were part of what the sentence says, the sentence would have no determinate content in such contexts, and speakers who differ on the completion would understand it differently. If the restriction is pragmatic, indeterminacy is expected: what a speaker communicates is often less determinate than what her sentence says, and communication succeeds as long as the candidate completions agree on the object. The Russellian content is determinate; the communicated content need not be.

\paragraph{Explicit contrast.} A speaker who has just said \enquote{the table is covered with books} can go on, without correcting herself, \enquote{the table in the other room is bare.} If the first description had the semantic content \enquote{the table in this room}, the second would contrast with it as one restricted description with another. On the pragmatic account the second sentence states explicitly the restriction that the first left to the hearer, and its explicitness is what lets it mark a contrast. The ease with which speakers make implicit restrictions explicit, and override them, is what one expects of pragmatic material.

\section{Objections}

\begin{objection}[Varying restrictions]
In \enquote{the table is bigger than every other table}, said in a furniture shop, the description is restricted to the table at hand and the quantifier to the tables in the shop. Two restrictions in one sentence must be computed in the semantics, since a pragmatic process applies to whole utterances \citep{StanleySzabo2000}.
\end{objection}

\begin{reply}
Pragmatic enrichment is not confined to adding conjuncts to whole utterances. A hearer may enrich a constituent, as when \enquote{she took out her key and opened the door} is understood with the key opening the door. Principle P2 of LNG-19-a constrains only the source of the enriched material, the hearer's reasoning about the speaker's purposes, not its position. Two restrictions in one sentence are two enrichments.
\end{reply}

\begin{objection}[Attitude reports]
\enquote{John believes that the table is covered with books} reports a belief about the table in the room. If the embedded description had its Russellian content, the report would ascribe to John a belief that there is exactly one table in the universe, which he does not have.
\end{objection}

\begin{reply}
The report is literally false on that reading, and the speaker communicates the true report that John believes, of the table in the room, that it is covered with books. The pattern is the same as in unembedded uses, and it is the same as for quantifiers in attitude reports: \enquote{John believes that everyone came} does not ascribe to John a belief about everyone in the universe. A semantic theory that restricted the embedded description would need the variable to be valued by the speaker's context in some reports and by John's in others; the pragmatic account needs only that the hearer ask what the speaker means. LNG-19-a answers the general argument from truth-value intuitions (its Objection 2); the embedded case is an instance.
\end{reply}

\section{Conclusion}

Incomplete descriptions would be used as we use them in a language stipulated to be Russellian; their indeterminate completions are what pragmatic enrichment predicts and semantic restriction cannot accommodate; and restrictions that vary within a sentence or occur under attitude verbs are instances of local pragmatic enrichment. Domain restriction for definite descriptions is pragmatic, and definite descriptions are not context-sensitive.

\begin{thebibliography}{9}\small
\bibitem[Kripke(1977)]{Kripke1977} Kripke, S. (1977). Speaker's reference and semantic reference. \emph{Midwest Studies in Philosophy}, 2, 255--276.
\bibitem[Stanley and Szab\'o(2000)]{StanleySzabo2000} Stanley, J. and Szab\'o, Z.~G. (2000). On quantifier domain restriction. \emph{Mind \& Language}, 15(2--3), 219--261.
\bibitem[Strawson(1950)]{Strawson1950} Strawson, P.~F. (1950). On referring. \emph{Mind}, 59(235), 320--344.
\bibitem[Wettstein(1981)]{Wettstein1981} Wettstein, H.~K. (1981). Demonstrative reference and definite descriptions. \emph{Philosophical Studies}, 40(2), 241--257.
\end{thebibliography}

\end{document}
